\scanpage{38}
% archive foliation: 37
Soient $X_i$ sur $Y$ ($1\leq i\leq n$), des $F_i$ sur les $X_i$, \deleted{\unsure{on a}} un changement de base $Y'-Y$, et on compare
\[
\underline{\operatorname{Tor}}_*^{\underline O_Y}(F_1,-F_n)\otimes_Y Y'\qquad\text{et}
\]
\[
\underline{\operatorname{Tor}}_*^{\underline O_{Y'}}(F'_1,-F'_n)
\]
\margin{\diagram{\begin{tikzcd}[column sep=small,row sep=large]
X_1&X_2&X'_1&X'_2\\
Y&&Y'&\\
S&&S'&
\end{tikzcd}}}

On sait qu'on a un hom canonique du premier ds le second. À quelles conditions est-ce un isomorphisme ?

\deleted{Comme le dernier \unsure{est un} \unsure{l'aboutissement}}
\deleted{\unsure{D'une}} \inserted{On a des} suites spectrales \inserted{\unsure{ab. = le dernier,|ab. le dernier,}} commençant par
\[
\operatorname{Tor}_p^{\underline O_S}\bigl(\operatorname{Tor}_q^{\underline O_Y}(F_1,-,F_n)\deleted{\otimes_Y},\underline O_{S'}\bigr),
\]
il suffit que les $\operatorname{Tor}_q^{\underline O_Y}(F_1,-,F_n)$ soient $Y$-plats.
